\section{The Second Precept of Love}

\begin{quotationx}
This passage is taken from \emph{Dio e il Poeta} by \textbf{Guido de Giorgio}. \textbf{Julius Evola} asserted that De Giorgio followed a Vedantized Christianity; this is immediately apparent in his description of the world as the ``play of God'', which has its exact analogue in the idea of \emph{lila\footnote{\url{https://en.wikipedia.org/wiki/Lila\_(Hinduism)}} in Hinduism}. Readers will also recognize the influence of \textbf{Rene Guenon}`s \emph{Multiple States of Being}, as shown in the topics of the Infinite, Possibility, or multiple worlds. Then again, Guenon claimed only to be explicating the metaphysics of the \textbf{Vedanta}. Yet, De Giorgio does not approach that text as a logical exercise, but rather as ``one who knows''.

Finally, there is the profound influence of \textbf{Saint Thomas Aquinas} and, of course, \textbf{Dante}. Much of the vocabulary, which cannot come across in a translation, is archaic, originating from Dante's time. There are several references to \emph{The Divine Comedy}, which I have provided whenever I could recognize them. The Comedy is more than a poem, it is a detailed text for initiates. It seems to us that De Giorgio followed the path outlined, climbing the seven story mountain of Purgatory, and saw the world from that height.
\end{quotationx}


It is tied to the first, and subordinated to the first, \emph{ancillarily}, but in the purely religious meaning in which this term must be used, as profound comprehension of the divine drama, of creation, that, divine and nonhuman, refers to that which, coming from God, returns to God, as a ``restitution'' which becomes a ``gift'', a gift of love. The world is the fruit and the work of divine love which multiplies itself, remaining one, gives itself because given to Him, offers itself because offered to Him, breaks itself, shatters itself, apparently because every part, every fraction, purely apparent, comes back in the All, from which it separated itself as pure ``play''. The world is God's ``play'', removing from this word any frivolous, human limitation whatsoever, in the special, precise sense that, e.g., Saint Thomas gives it, when he speaks of contemplation, in the Latin word ``ludus''. The world is the ``mirror'' of God, the crystal in which the Unity contemplates Himself and, contemplating Himself , pluralizes Himself , multiplies Himself , still remaining One, because He cannot not be One while being One, like a fire of a thousand sparks each one of which is fire but not \emph{the fire}, since every relation with God is \emph{uni}lateral, not double, not reciprocal, God alone being God and what has life from God, has being, not being able to be God, if God is God, ``for the contradiction does not allow it''\footnote{\emph{Inferno}, XXVII}. The world is ``imprinted'' by God in the sense that God, while passing, creates, and the created reflects Him, it is His footprint, nor can it be considered otherwise than in relation to Him, without which nothing would be, it is obvious, because the trace is a sign, a sign that indicates the Giver, the Donor, the One who imprinted it, stamped it, engraved it. One could not speak of God without the world, although God is without the world, beyond all possible worlds, for the simple reason that He cannot exhaust Himself in creation which is His work, nor limit Himself to it, why His power is always ``in infinite excess''\footnote{\emph{Paradise}, XIX} in the ambit of created things and since this ``excess'' is infinite, it is understood without difficulty, that which we call ``everything'' created, is, ontologically, nothing in the face of He who created it.

If we call the ``world'' the visible, it is clear that is cannot reflect God who, by definition, by necessity, by absoluteness, is invisible, that not removing, however, what is visible, insofar as it is visible, it is also the sign, imprint, mirror, reflection of the invisible, and that the visible breaks itself, shatters itself indefinitely (a multitude of beings) \emph{because it cannot contain the invisible}. The creaturely multiplicity is the indication of this impotency, as numerical indefiniteness is the indication of the infinity of unity, without which it could not offer itself, although no number of the series entirely reflects the unity that is in everything and in none of them.

God as God must create the world, without which He would not be God, or, at least, He would not be known and therefore loved as God, but if God is the Invisible, it is also necessary that he is God outside and beyond the world, outside and beyond creation, which is one of his possibilities but not all of his Possibility that must be infinite as He himself is infinite. His ``infinite excess''... Manifestly, therefore, everything originating from Him, each thing should return to Him, and in fact, does return to him, must return to Him, but each thing, as every creature, through its way and manner, or \emph{diversimode} diverse modes, Latin as Saint Thomas would say, by levels of integration. Among all creatures, there is man, who has his mystery which corresponds to the Mystery of the Incarnation, the descent for the ascent, and this mystery can determined in a certain manner saying that man can ``take up again'' all the work of God, in a \emph{syn-thesis}, i.e., in a \emph{com-position}, in the universality of understanding, of prehension as the eye which embraces an entire horizon in a unitary vision, uniting the multiplicity, joining the separated, but it is necessary, however, that he has a \emph{third eye} in the middle of the forehead in which the two eyes are ``centered'', to resolve the duality in unity, since the two is limitation, it is a throwback, and must result in the three.

The privilege of man, therefore, is this, to be able to return to God, \emph{freely and totally} because he alone, among all creatures, has the option, the choice, to remain in creation, or return to the Creator, taking up all creation, \emph{capitulatim} summarily, Latin in the unity. This is the true, the great dignity of man that he owes to the faculty that, only in him, embraces not only what appears (the visible, the world), but that which is, God. This faculty is the intellect, not reason which, visibly, cannot pass beyond the visible domain of quantity, the world, creation.


\flrightit{Posted on 2023-03-21 by Cologero}

\begin{center}* * *\end{center}

\begin{footnotesize}
\begin{sffamily}
\texttt{William Zeitler on 2023-03-22 at 09:32 said:}

What's the first precept of love? I haven't been able to find a blog by that title or a link thereto. Thanks!


\hfill

\end{sffamily}
\end{footnotesize}

